% Transcription — fonds Alexandre Grothendieck, shelfmark 82, batch 3 (pages 41-60).
% Established from the facsimile archives/batches/82/batch-03.pdf, cut from a
% copy of 82.pdf fetched through the project's relay
% (https://grothendieck-relay.onrender.com/source/82.pdf); PDF page k =
% archivists' page 40+k, no cover sheet. Per the skill transcribe-grothendieck.
% The folder is seventy-eight pages in four batches (1-20, 21-40, 41-60,
% 61-78), transcribed in parallel by four passes; neither neighbouring file
% existed when this one was written, so nothing here was checked against them.
% Pass: Opus 5.5 (claude-opus-5-5), 2026-09-26 — first pass, unchecked against the pages by a human.
%
% Dating and title. \dating is the folder's own catalogue date, copied
% verbatim; the folder belongs to the group « Géométrie et topologie
% combinatoire » (69-89), but since it has a date of its own the group's range
% is not added. \foldertitle is the catalogue title without its final full stop.
%
% Pages skipped: 45, a blank cream sheet (only show-through of p. 44).
% Page 56 is a pink cover inscribed in pencil « Dictionnaire lié à
% GP(1) ≃ SO(3) »; by the cover rule it takes no \page marker and its words
% become the \section* heading before p. 57, with a \note.
% Pages summarised: none. Pages transcribed from his hand: 41-44, 46-55,
% 57-60, all in blue felt-tip; p. 44 carries one line only, p. 55 half a page.
% Drawings: none in this batch.
%
% His pagination: none on the leaves. His own numbered runs: equations
% (1)-(25) on pp. 46-54 (one continuous run, from (1) M = V̌ ⊗ V to (25)
% T = V̌(A)*/G_m); a fresh run on p. 60, (1) and (2), inside « II Opérations
% de G_0 ». Enumerated items: a), b), c) on p. 42 (c) cancelled), a)-d) on
% p. 51, a), b), c) on p. 50 (cancelled); in the Dictionnaire, the list of
% categories C_0 ... C_5, C'_i, C''_i on p. 57 and « I Représentants
% typiques » 1)-6) on pp. 59-60, then « II Opérations de G_0 » 1)-4) on p. 60.
% On p. 59 he counts « 11 catégories » while p. 58 counts 9 × 8 = 72
% equivalences (nine once C''_1, C''_2 are struck); both are left as written.
% The margin of p. 55 refers, doubtfully, to « p. 52 »; kept as plain text.
%
% What the batch is about. pp. 41-44: the end of a proof, begun in batch 2,
% on a finite étale subgroup μ of prime order p (odd) in a form G of GP(1):
% its centraliser H is smooth commutative of relative dimension 1, and its
% normaliser N has N° = H, N/H ≃ {±1}_S acting by h ↦ h^{-1} (p. 42 states it
% and starts the proof via SL(V)/μ_2), then a question whether every smooth
% subgroup of relative dimension 1 is the centraliser of its Lie algebra
% (p. 43), which breaks off. pp. 46-55: quaternion algebras M = End(V),
% V of rank 2: the equivalences M ↦ G, G̃, (E, q, ω), (P ⊃ X), X; the
% isomorphisms X ≃ P(V), Q = Var(Δ) ≃ P(V) × P(V) via u ↦ (Ker u, Im u);
% the involution u' = Tr(u)·1 − u; special commutative subalgebras A of rank
% 2, the proposition (p. 51) on lines L stable under A, P̌ ≃ Div_2^+(X), Y ≃
% Spec A, and the « jacobian » subgroups T = V̌(A)*/G_m of G (smooth
% connected commutative of relative dimension 1). pp. 57-60, under the cover
% « Dictionnaire lié à GP(1) ≃ SO(3) »: the plan of a dictionary of
% equivalent fibred categories (gerbes) over a scheme S — torsors under
% GP(1), forms of GP(1), SL(2), gp(1), sl(2), quaternion algebras, conic
% bundles, pairs (P, X), special quadratic modules of rank 3 — the typical
% representatives over Spec Z (q_0 = xy + z^2, M_0 = M_2(O)), and the start of
% the actions of G_0 on them.
%
% Hardest pages: 41 and 51-54, fast cursive whose connective prose often does
% not carry; 57 and 59 are lists overwritten with struck lines. Notation
% fixed for the batch: GP(1) as \mathrm{GP}(1); his round capital for the
% jacobian subgroup as \mathcal{T}; the gothic letter on pp. 58-59 as
% \mathfrak{X}; the centre on p. 60 as \mathfrak{z}.

\documentclass[11pt,a4paper]{article}
% Le préambule commun à toutes les transcriptions du fonds Grothendieck.
%
% Il tient en une idée : **l'incertitude se compose, elle ne se résout pas.**
% Une transcription qui lisse un mot douteux a détruit la seule chose pour
% laquelle on la faisait — un lecteur doit pouvoir savoir, à chaque endroit,
% ce qui était sur le feuillet et ce qui a été ajouté. Les six macros qui
% suivent sont donc l'appareil critique, et non de la décoration.
%
% Les mêmes commandes sont reconnues par `scripts/render.mjs`, qui produit la
% vue de lecture du site. Une macro ajoutée ici doit l'être aussi là-bas,
% faute de quoi le rendu échouera bruyamment — ce qui est voulu : un rendu
% silencieusement faux est pire qu'un rendu absent.

\ProvidesPackage{grothendieck}[2026/08/08 Transcriptions du fonds Grothendieck]

\usepackage{amsmath,amssymb,amsthm}
\usepackage{mathrsfs}
\usepackage{stmaryrd}
% Les diagrammes commutatifs sont partout dans ce fonds : c'est la langue
% même de la géométrie algébrique de Grothendieck, et une transcription qui
% les rendrait en prose ne transcrirait rien.
\usepackage{tikz-cd}
% Français par défaut — c'est la langue des feuillets — mais l'anglais est
% chargé aussi : la traduction et le résumé sont des documents anglais, et la
% typographie française (espaces avant les deux-points) y serait une faute.
% Ces documents basculent par \selectlanguage{english} en tête de corps.
\usepackage[english,french]{babel}
\usepackage{fontspec}
\usepackage{geometry}
\geometry{a4paper,margin=25mm,marginparwidth=28mm}
\usepackage{marginnote}
\usepackage{xcolor}
% ulem, not soul. `soul`'s \st and \ul are built on a character-by-character
% scan that XeLaTeX's font machinery breaks: the document fails with an
% undefined control sequence rather than degrading. ulem does the same two jobs
% and is Unicode-engine safe. `normalem` stops it hijacking \emph, which the
% transcriptions use for Grothendieck's own emphasis.
\usepackage[normalem]{ulem}
\usepackage{hyperref}
\hypersetup{colorlinks=true,linkcolor=black,urlcolor=black,pdfborder={0 0 0}}

% --- Métadonnées du lot -----------------------------------------------------
% Reprises telles quelles par le script de rendu, qui les affiche en tête de
% la vue de lecture. Elles fixent ce que le fichier prétend couvrir : sans
% elles, une transcription flotte sans point d'attache dans un fonds de seize
% mille feuillets.

\newcommand{\folder}[1]{\gdef\grothFolder{#1}}
\newcommand{\batch}[1]{\gdef\grothBatch{#1}}
\newcommand{\leaves}[2]{\gdef\grothFirst{#1}\gdef\grothLast{#2}}
\newcommand{\foldertitle}[1]{\gdef\grothFoldertitle{#1}}
\gdef\grothFolder{?}\gdef\grothBatch{?}\gdef\grothFirst{?}\gdef\grothLast{?}\gdef\grothFoldertitle{}

% --- Le résumé --------------------------------------------------------------
%
% Il ouvre la lecture modernisée, sous le titre « Résumé », et il est composé
% en petit. Ce n'est pas un ornement : c'est le seul passage du document qui
% ne soit pas des mathématiques, et un lecteur doit pouvoir le reconnaître
% avant de l'avoir lu — puis le sauter s'il connaît déjà le sujet, ou s'y
% arrêter s'il ne le connaît pas. Le filet qui le suit marque la reprise du
% corps.

\newenvironment{resume}
  {\small\setlength{\parskip}{0.45em}}
  {\par\medskip\hrule\medskip}

% --- Mots-clés ---------------------------------------------------------------
%
% En anglais, en clôture du résumé de la lecture modernisée : le vocabulaire
% moderne sous lequel chercher ce que les feuillets construisent. Le manifeste
% du site (`npm run manifest`) les extrait pour étiqueter la cote entière —
% cette ligne est la source unique des tags, il n'y a pas de fichier à côté.
\newcommand{\keywords}[1]{\par\smallskip\noindent
  {\footnotesize\textsc{Keywords} --- \textit{#1}}\par}

% --- L'appareil critique ----------------------------------------------------

% Le numéro de feuillet, en marge. C'est l'ancre qui permet de régler un
% désaccord sur une formule en regardant *un* feuillet plutôt qu'une chemise
% de six cents. Le site s'en sert aussi pour tourner les pages du fac-similé
% quand on fait défiler la transcription.
\newcommand{\leaf}[1]{%
  \marginnote{\footnotesize\color{gray!70}\textsf{#1}}%
  \label{leaf:#1}\ignorespaces}

% Illisible : on ne devine pas. Un mot inventé qui se lit comme les autres est
% le pire résultat possible.
\newcommand{\ill}{\textcolor{red!60!black}{[\,\ldots\,]}}

% Lecture douteuse : le mot est donné, et signalé comme incertain.
\newcommand{\uncertain}[1]{\uline{#1}}

% Ajout de l'éditeur — une lettre manquante, un mot sous-entendu.
\newcommand{\add}[1]{\textcolor{blue!50!black}{[#1]}}

% Biffé par Grothendieck lui-même. Ce qu'il a rayé fait partie du document :
% une rature dit souvent où la pensée a changé de direction.
\newcommand{\struck}[1]{\sout{#1}}

% Note du transcripteur, sur l'état matériel du feuillet.
\newcommand{\note}[1]{\par\smallskip{\small\color{gray!60!black}\textit{[#1]}}\par\smallskip}

% Note marginale de Grothendieck, distincte du corps du texte.
\newcommand{\marginal}[1]{\par\smallskip\noindent\hspace{1em}%
  {\small\color{gray!60!black}#1}\par\smallskip}

% --- Titre ------------------------------------------------------------------

\AtBeginDocument{%
  \begin{center}
    {\large\bfseries Fonds Alexandre Grothendieck --- cote n\textsuperscript{o}~\grothFolder}\\[2pt]
    {\itshape\grothFoldertitle}\\[4pt]
    {\small feuillets \grothFirst--\grothLast\ (lot \grothBatch)}\\[2pt]
    {\footnotesize\color{gray!60!black}Transcription établie d'après le fac-similé numérisé
     par l'Université de Montpellier.\\
     Les lectures incertaines sont soulignées, les illisibles notées [\,\ldots\,],
     les ajouts entre crochets.}
  \end{center}
  \vspace{1em}}

\endinput


\folder{82}
\batch{3}
\pages{41}{60}
\dating{[vers 1982]}
\watermark{Édition de démonstration}
\foldertitle{SO(3) ≃ GP(1) : notes manuscrites (s.d.)}

\begin{document}

\page{41}
\note{la page commence au milieu d'une démonstration, commencée sur la page précédente (batch 2).}
mais il est facile de le \uncertain{vérifier}, en choisissant un
\uncertain{générateur} de $\mu$, et en introduisant les valeurs propres
$\zeta, \zeta^{-1}$ de $u$ \ill{} $\lbrace \zeta^{p} = 1$
\marginal{\ill{} \uncertain{racines} de $p \geq 2$, \uncertain{primitives}
\ill{}}
\note{insertion interlinéaire au-dessus de la ligne, en partie soulignée ; lecture très incertaine.}
donc $\zeta \neq \zeta^{-1}$ car $p$ \uncertain{impair} $p \neq \pm 1\rbrace$
\uncertain{puisque} \ill{} \uncertain{fid.} de $\mathbf{Z}/p\mathbf{Z}$) le
\uncertain{commutant} \add{de $u$} est formé des $(\lambda, \lambda^{-1})$, le
\uncertain{normalisateur} est formé des matrices qui sont ou de cette forme
diagonale, ou le produit d'une telle par une \uncertain{symétrie} \ldots\ D'où
\[
  N/H \simeq \lbrace \pm 1 \rbrace_{S} \subset (\mathbf{Z}/p\mathbf{Z})^{\times}_{S} .
\]
\note{après $(\mathbf{Z}/p\mathbf{Z})^{\times}$ un petit signe raturé.}

Si $p = \operatorname{car} k$, \uncertain{alors} $u$ est de la forme
\[
  u = \mathrm{id}_V + v
\]
avec \struck{\ill{}}
\[
  v^2 = 0, \quad v \neq 0 ,
\]
et le commutant est formé des matrices \emph{\uncertain{unipotentes}} de la
forme $\alpha\,\mathrm{id} + \beta v$ \struck{il faut \ill{} et $\alpha^2 = 1$
donc \ill{}} la condition \uncertain{unipotente} \emph{signifie} $\alpha^2 = 1$
(condition \ill{} \add{$\alpha = \pm 1$}, puisque $p \neq 2$ !)
\note{« $\alpha = \pm 1$ » est inséré au-dessus de la ligne ; le signe entre $\alpha$ et $\pm 1$ peut être aussi $\neq$.}
et on a encore la lissité, la connexité de $U$, etc.
\struck{\ill{} fait que \ill{}}

b) \emph{Cas général} (il s'agit de prouver que $N$ est lisse et $N^{\circ}$
commutatif) -- Comme \struck{$N/N^{\circ}$ est un groupe étale \ill{}}
l'image inverse par $N \to \underline{\mathrm{Aut}}(\mu) = (\mathbf{Z}/p\mathbf{Z})^{\times}_S$
de la section unité est un sous-groupe \uncertain{ouvert} de $N$ qui coïncide
avec $N^{\circ}$ en chaque fibre, $N$ \ill{} \uncertain{égal} : $N^{\circ}$ ; de plus
$N/N^{\circ} \subset (\mathbf{Z}/p\mathbf{Z})^{\times}_S$ \ill{}
\note{la page s'arrête ici ; la suite n'est pas à la page 42, qui commence un nouvel énoncé.}

\page{42}
Soit $G$ une forme de $\mathrm{GP}(1)$, \struck{\ill{}} $p$ un nb premier,
impair, et $\mu$ un sous-groupe fini étale de $G$ de rang $p$. Je dis que
\begin{itemize}
  \item[a)] Le centralisateur \add{$H$} de $\mu$ est lisse (de dim relative 1
  commutatif), à fibres connexes \uncertain{hors} car. \uncertain{rés.} $\neq p$.
  \note{« $H$ » et « commutatif » sont insérés au-dessus de la ligne ; « à fibres connexes \ldots{} $\neq p$ » est écrit en interligne sous la ligne.}
  \item[b)] Soit \struck{\ill{}} $N$ le normalisateur de $\mu$. Alors $H = N^{\circ}$ ;
  $N/H \simeq \lbrace \pm 1\rbrace_S$ et opère sur $H$ par
  \struck{\ill{}} $h \mapsto h^{-1}$)
  \struck{donc $N/H$ \ill{} $\underline{\mathrm{Aut}}_{\mathrm{gr}}(\mu) \simeq (\mathbf{Z}/p\mathbf{Z})^{\times}_S = \mathbf{F}^{\times}_{p\,S}$}
  \note{la ligne « $N/H \simeq \lbrace\pm1\rbrace_S$ et opère sur $H$ par \ldots{} $h \mapsto h^{-1}$) » est écrite en interligne au-dessus de la ligne biffée.}
\end{itemize}

\struck{c) Il existe une section unique $\sigma$ de $N/H$ dont l'image dans
$\Gamma(S, (\mathbf{F}_p^{\times})_S)$ est la section constante de valeur $-1$.
Son image par $N/H \hookrightarrow \underline{\mathrm{Aut}}_{\mathrm{gr}}(H)$ est
également la symétrie $h \mapsto h^{-1}$.}
\note{tout ce paragraphe est encadré et barré de grands traits en diagonale ; devant « Il existe » un début raturé ; dans « l'image dans \ldots{} » les mots $\underline{\mathrm{Aut}}$ et $\mathbf{Z}/p\mathbf{Z}$ sont biffés au profit de $\Gamma(S, (\mathbf{F}_p^{\times})_S)$.}

\emph{Dém.} OPS $G = \overbrace{\mathrm{SL}(\underline{V})}^{\widetilde{G}}/\mu_2$,
$\underline{V}$ fibré vectoriel. Comme $p$ impair, on \struck{peut}
remonter $\mu$ en $\widetilde{\mu} \simeq \mu$ dans $\widetilde{G}$, de façon unique.
Soit $\widetilde{H}$ son centralisateur dans $\widetilde{G}$, il \uncertain{suffit} de
\uncertain{voir} que $H = \widetilde{H}/\mu_2$ est \uncertain{son} centralisateur
dans $G$. On a \struck{\ill{}} prouvé que $H$ est lisse commutatif de dim relative 1.

a) $S = \operatorname{Spec} k$, $k$ \add{corps} alg. clos. Si $p$ premier à la
car, le résultat \add{de lissité} est connu pour \uncertain{tt} \uncertain{sous}-groupe de
\uncertain{type} mult. dans un groupe lisse,

\page{43}
Tout sous-groupe lisse \add{\uncertain{rel.} commutatif} \struck{de rang} \add{de}
dim rel.~1 dans $G$ est-il de la forme $C(\mathfrak{h})$ -- i.e.\ est-il
déterminé par Lie $H$ ? Oui \add{s'il est commutatif et} si les car.\
résiduelles sont $\neq 2$, car il est contenu alors dans $C(\mathfrak{h})$ qui a
les \uncertain{mêmes} \uncertain{fibres}, donc lui est égal. \uncertain{Mais} on se
demande si la condition de commutativité est automatiquement satisfaite ?
\note{« rel.\ commutatif » et « s'il est commutatif et » sont des insertions interlinéaires ; « de rang » est barré d'un trait ondulé.}

\emph{Question \uncertain{infinitésimale}} \add{\ill{}}. On sait que c'est le cas
sur un corps -- tt groupe lisse \add{connexe} de dim 1 est commutatif -- donc la
question est une question de déformations de sous-groupes, liée à un $H^2$ de
Hochschild.

Considérons le normalisateur de $H$ -- qui est aussi celui de $\mathfrak{h}$,
puisque la corr.\ entre $H$ et $\mathfrak{h}$ est 1-1. J'ai envie de prouver
qu'il est (dans le cas car.\ rés.\ $\neq 2$) lisse de dim relative 1 -- mais
c'est faux \struck{\ill{}} sur un corps, quand $H$ est unipotent, le normalisateur
devient un Borel ! C'est ici qu'il faut \uncertain{finalement}

\page{44}
une structure décrite par \emph{\uncertain{rigidifiée}} \ldots\ \struck{\ill{}}
\note{le reste de la page est blanc ; la phrase commencée en bas de la page 43 s'arrête ici.}

\page{46}
À l'Algèbre de quaternions \add{$\mathcal{M}$}, \uncertain{on} associe
\begin{itemize}
  \item forme $G$ de $\mathrm{GP}(1)_S$ ;
  \item -- $\widetilde{G}$ -- $\mathrm{SL}(2)_S$ ;
  \item -- $\mathfrak{g}$ de $\mathfrak{gp}(1)_S$ ;
  \item -- $\widetilde{\mathfrak{g}}$ -- $\ill{}(1)_S$ ;
  \item $(E, q, \omega)$ ($E = \widetilde{\mathfrak{g}}$) \uncertain{module} spécial quadratique de rg~\uncertain{2} ;
  \item $P \supset X$ quadrique dans plan projectif ($P = \check{\mathbb{P}}(E)$, $X = \mathrm{Var}(q)$) ;
  \item $X$ fibré en droites projectives.
\end{itemize}
\note{la liste est tenue par une grande accolade ; les tirets répètent « forme \ldots{} de ». Les deux premières lignes ($G$, $\widetilde{G}$) et les deux suivantes ($\mathfrak{g}$, $\widetilde{\mathfrak{g}}$) sont regroupées par de petites accolades, avec une flèche verticale entre les deux groupes. Le modèle de $\widetilde{\mathfrak{g}}$ est surchargé et illisible ; on attend $\mathfrak{sl}(2)_S$ (cf.\ page 59). « rg 2 » est la lecture de l'encre, alors que $E = \widetilde{\mathfrak{g}}$ est de rang 3.}

On trouve que
\[
  \mathcal{M} \mapsto G, \quad \mathcal{M} \mapsto \widetilde{G}, \quad
  \mathcal{M} \mapsto (E, q, \omega), \quad
  \mathcal{M} \mapsto (P \supset X), \quad \mathcal{M} \mapsto X
\]
sont des équivalences de catégories \add{fibrées}, et que
$\mathcal{M} \mapsto \mathfrak{g}$, $\mathcal{M} \mapsto \widetilde{\mathfrak{g}}$
sont des équiv.\ de cat.\ fibrées au-dessus de $\operatorname{Spec} \mathbf{Z}(\tfrac{1}{2})$.

Considérons maintenant le cas où \ill{}
un fibré vect.\ de rang 2 $V$, et \struck{\ill{}}
\[
  (1) \qquad \mathcal{M} = \check{V} \otimes V \simeq \underline{\mathrm{End}}(V)
\]
On va d'abord définir un iso can.
\[
  (2) \qquad X \overset{\varphi}{\simeq} \mathbb{P}(V) \;\bigl(\simeq \mathbb{P}(\check{V})\bigr)
\]
\note{« orth. » est écrit au-dessus du second $\simeq$, et relié par un trait à l'encadré de la ligne suivante.}
Pour cela, si $L \subset V$ \struck{\ill{}} est un sous-fibré vectoriel inv.\ de $V$,
\struck{\ill{}} $L' \subset V'$ son orthogonal ($L' = L^{\perp}$), on lui associe
\[
  (3) \qquad \varphi(L) = L' \otimes L = L^{\perp} \otimes L \subset \check{V} \otimes V = \mathcal{M}
\]
et on trouve bien un iso \ldots

\page{47}
Soit
\[
  (4) \qquad Q = \mathrm{Var}(\Delta) \subset \check{\mathbb{P}}(\mathcal{M})
\]
où $\Delta$ est la forme quadratique
\[
  (5) \qquad \Delta(u) = -\det u
\]
sur $\mathcal{M}$ -- qui est non dégénérée. On va définir un isomorphisme canonique
\[
  (6) \qquad Q \xrightarrow[\sim]{\ \psi\ } \mathbb{P}(V) \times \mathbb{P}(V) \quad (\simeq X \times X)
\]
\note{avant $\mathbb{P}(V) \times \mathbb{P}(V)$, un premier « $X \times X$ » est biffé.}

En effet, $Q$ correspond aux \struck{$u \in \Gamma(\mathcal{M})$ \ill{}}
\add{$u : V \to V$} (non nuls sur chaque fibre) tels que $u$ soit
\emph{strictement de rang 1} i.e.\ se factorise canoniquement en
\[
  V \twoheadrightarrow I \hookrightarrow V
\]
\note{les deux flèches portent « épi » et « mono ».}
avec $I$ inversible, et il est connu \ill{} \add{\uncertain{multipl.}}.
Se donner \uncertain{u} \ill{} équivaut à se donner le quotient $I$ de $V$,
\struck{\ill{}} ou encore le noyau de $u$ ; et \uncertain{l'image} \ill{} $I$ de $V$
i.e.\ $\mathrm{Im}\, u$. Comme ceux-ci peuvent être pris indépendamment. D'où (6),
\uncertain{qui} est induit par
\[
  (7) \qquad u \longmapsto (\underbrace{\mathrm{Ker}\, u}_{L}, \underbrace{\mathrm{Im}\, u}_{M})
\]
\marginal{$u$ non nul sur ch.\ fibre, $\Delta(u) = 0$}
\note{écrit sous le $u$ de (7). La mise en mots de l'avant-dernière phrase (« Se donner \ldots{} Im $u$ ») est très incertaine ; les formules sont sûres.}
\marginal{\emph{NB} L'iso \uncertain{inverse} \ill{} \uncertain{encore} plus
\uncertain{simplement} \ill{} $(L, M) \mapsto L^{\perp} \otimes M \subset \check{V} \otimes V$}
\note{note écrite en diagonale à droite de (7).}

\page{48}
Soient \struck{\ill{}}
\[
  \underset{L^{\perp} \otimes L}{\underline{L}}, \ \underset{M^{\perp} \otimes M}{\underline{M}} \subset \mathcal{M} = \underline{\mathrm{End}}(V)
\]
correspondant à $L$, $M$, \struck{ils sont caractérisés} par la condition de
correspondre à des sections de $X$, i.e.
\[
  (8) \qquad
  \begin{cases}
    \underline{L} \subset \bar{\mathfrak{g}}, \ q \mid \underline{L} = 0 \\
    \underline{M} \subset \mathfrak{g}, \ q \mid \underline{M} = 0
  \end{cases}
\]
i.e.\ $\mathrm{Tr}, \det \mid \underline{L} = 0$, resp.\ $\mathrm{Tr}, \det \mid \underline{M} = 0$
\note{les deux « i.e. » sont écrits à droite de chacune des deux lignes de (8).}
\note{au-dessus de $\underline{L}$ et $\underline{M}$, reliés par des signes d'égalité verticaux, $L^{\perp} \otimes L$ et $M^{\perp} \otimes M$. Le premier $\mathfrak{g}$ de (8) porte un trait au-dessus, le second non ; c'est vraisemblablement le même objet ($\widetilde{\mathfrak{g}}$, les endomorphismes de trace nulle), noté de façon hâtive.}
et que de plus on ait, posant
\[
  (9) \qquad
  \begin{cases}
    J = \underline{O}_S \cdot u \subset \mathcal{M} & (J \in \Gamma(Q/S)) \\
    J \circ L = 0 \\
    \underline{M} \circ J = 0
  \end{cases}
\]
\note{un double trait horizontal sépare ce qui suit.}

Notons que la condition $J \in \Gamma(Q/S)$ implique que
\[
  (10) \qquad \mathcal{A} = \underline{O}_S \oplus J \hookrightarrow \mathcal{M} = \underline{\mathrm{End}}(V)
\]
\note{« unité » est écrit sous la flèche d'inclusion.}
est un \add{sous-}fibré vectoriel de rang 2, \struck{\ill{}} la donnée de $J$
\emph{équivaut} à la donnée \struck{\ill{}} de $\mathcal{A}$, \emph{et} d'un
homomorphisme d'Algèbres
\[
  (11) \qquad \mathcal{A} \xrightarrow{\ \rho\ } \underline{O}_S
\]
(savoir l'unique hom d'Alg.\ de noyau $J$).
\marginal{\uncertain{explicitement} \ill{}}
\note{note marginale à gauche, en face de (10)--(11), qu'un trait vertical rattache à ces lignes.}

Ainsi on a un iso

(12) \quad $Q \simeq$ ens des couples $\mathcal{A}, \rho$,

$\mathcal{A} \subset \mathcal{M}$ sous-Algèbre \struck{\ill{}} \uncertain{spéciale},
$\rho : \mathcal{A} \to \underline{O}_S$ hom d'Alg.

\page{49}
Mais il y a dans $Q$ une \uncertain{automorphie} \uncertain{involutive}
canonique, qui correspond à la symétrie des deux facteurs $X \times X$ dans (6).
Elle est induite par un anti-automorphisme \add{can.} de l'Algèbre $\mathcal{M}$
\struck{obtenu (d'ailleurs \ill{})} \struck{en choisissant une base}
\add{\ill{} base} ou de $\det \check{V}$ \add{$= \Lambda^2 \check{V}$}, d'où un isom
\struck{\ill{}}
\[
  (13) \qquad V \xrightarrow[\sim]{\ \alpha\ } \check{V} \quad \bigl(V \to \check{V} \otimes \det V \xrightarrow{\sim} \check{V}\bigr)
\]
\note{entre $\check{V}$ et la parenthèse, il écrit « composé de ».}
ce qui permet de définir
\[
  (14) \qquad u \longmapsto u' = \alpha^{-1}\, ({}^{t}u)\, \alpha
\]
\note{la ligne « obtenu \ldots{} base » est surchargée de biffures et d'insertions ; seules les formules sont sûres.}
Mais changer de $\alpha$ revient à multiplier $\alpha$ par une section unité
$\lambda \in \mathbb{G}_m(S)$, donc \add{$\alpha^{-1}$ devient $\lambda^{-1}\alpha^{-1}$}
\struck{\ill{}} et $u'$ ne change pas. Par construction
\[
  (15) \qquad (uv)' = v' u'
\]
En fait, la formule donne
\[
  (16) \qquad u' = \mathrm{Tr}\, u \cdot 1 - u
\]
plus précisément \ldots
\[
  (17) \qquad
  \begin{cases}
    u + u' = (\mathrm{Tr}\, u)\, 1 \\
    u u' = (\det u)\, 1
  \end{cases}
\]
et il est immédiat que $u \mapsto u'$ échange Ker et Im pour les $u$ tels que
$\Delta(u) = 0$, $u$ non nul sur les fibres \ldots
\note{« $\Delta(u) = 0$ » : l'encre donne plutôt $\Gamma(u) = 0$ ; la forme quadratique de (5) est $\Delta$.}

Quand on voit ce que devient cette involution de $\mathcal{M}$ et $Q$ en termes de
(12), on voit que : le couple $(\mathcal{A}, \rho)$ \uncertain{lui} associe
$\mathcal{A}$ ($= \mathcal{A}'$ !) et $\rho'$ défini par
\[
  (18) \qquad \rho'(u) = \rho(u')
\]
\marginal{(\emph{NB} $u \mapsto u'$ est un \uncertain{automorphisme} \ill{}, car $\mathcal{A}$ est commutatif)}

\page{50}
Considérons l'hom
\[
  (19) \qquad Q \longrightarrow \check{\mathbb{P}}(\mathfrak{g}) = \mathbb{P}(\widetilde{\mathfrak{g}}) \quad \bigl(= \Sigma\bigr)
\]
\note{la parenthèse porte : « = schéma des s-Algèbres spéciales de $\mathcal{M}$ », avec $\Sigma$ écrit dessous et relié par un signe d'égalité vertical.}
C'est évidemment un \emph{épi}, et si $\mathcal{A}_{\Sigma}$ est « la
sous-Algèbre spéc.\ universelle » de $\mathcal{M}_{\Sigma}$ -- on trouve donc un iso.\ can.
\[
  (20) \qquad Q \simeq \operatorname{Spec}(\mathcal{A}_{\Sigma})
\]
\note{un double trait horizontal sépare ce qui suit.}

Ceci étant, $\mathcal{A}$ étant fixée, et $\underline{L}, \underline{M}$ des
sous-fibrés inv.\ de $\widetilde{\mathfrak{g}}$ correspondant à des sections de $X$
(i.e.\ sur lesquels \uncertain{$q$} s'annule) \add{$= L^{\perp} \otimes M$},
à quelle condition le $J \subset \mathcal{M}$ section de $Q$ qui leur
correspond est-elle $\subset \mathcal{A}$ ? (Les couples $\underline{L}, \underline{M}$
en question correspondent donc exactement aux sections de $Y = (\ldots)$.
\note{la fin de la phrase est « sections de $Y$ = (pour $\mathcal{A}$ sur $S$.) » ; l'insertion « $= L^{\perp} \otimes M$ » est écrite au-dessus de $J$.}

\struck{Je dis qu'il \ill{} et \uncertain{suffit} que $\underline{L} \subset \mathcal{A}^{\perp}$}

\struck{Les conditions équivalentes, pour $\mathcal{A}, L, M$ donnés,
$\mathcal{A}$ s-Alg.\ spéciale de $\mathcal{M}$, $L, M$ ss-fibrés
\add{de rang 1 de $V$} :}

\struck{a) $\underline{L} \subset \mathcal{A}^{\perp}$ (i.e.\ $\mathrm{Tr}\, \underline{L} J = 0$ si $\mathcal{A} = \underline{O}_S \oplus J$)}

\struck{b) $\underline{M} \subset \mathcal{A}^{\perp}$ (i.e.\ $\mathrm{Tr}\, \underline{M} J = 0$ \ldots)}

\struck{c) $L^{\perp} \otimes M \subset \mathcal{A}$}

\struck{De plus, pour $\mathcal{A}$ \add{et} $L$ \add{$\in X(S)$} satisfaisant
\add{$L \subset \mathcal{A}^{\perp}$}, il existe un et un seul $M \in X(S)$
satisfaisant c) [donc}
\note{tout ce bloc est barré de deux grands traits obliques et d'un trait horizontal en tête ; « satisfaisant $\mathrm{Tr}, \det \mid L, M = 0$ » y est lui-même biffé. Les étiquettes « a) » « b) » sont surchargées et se lisent mal ; la page s'arrête sur « [donc ».}

\page{51}
\emph{Proposition.} Soient $\mathcal{A}$ une ss-Alg.\ spéciale, et $L$ un
sous-fibré \uncertain{droite} de $V$, $\underline{L} = L^{\perp} \otimes L \subset \mathcal{M}$
le sous-fibré isotrope de $\widetilde{\mathfrak{g}}$ associé. Conditions équivalentes
\begin{itemize}
  \item[a)] $\underline{L} \subset \mathcal{A}^{\perp}$
  \item[b)] $\mathcal{A}$ invarie la droite $L$ (\emph{NB} si $\mathcal{A} = \underline{O}_S + J$, \add{cela} revient à \uncertain{dire} que $J \cdot L \subset L$)
  \item[c)] $\exists$ \add{\uncertain{et unique}} ss-fibré inv.\ $J$ de $\mathcal{A}$ tel que $J \cdot L = 0$
  \item[d)] $\exists\, M \subset V$ ss fibré vectoriel \uncertain{droite} tel que $L^{\perp} \otimes M \subset \mathcal{A}$
\end{itemize}
\note{la proposition est marquée d'un trait vertical dans la marge gauche ; « d) » est précédé d'un signe raturé.}

De plus, \struck{\ill{}} sous ces conditions, le $M$ dans c) est unique, et
l'unique hom $\mathcal{A} \to \underline{O}_S$ de noyau $L^{\perp} \otimes M$
\struck{\ill{}} est aussi l'hom $\mathcal{A} \xrightarrow{u} \underline{O}_S \simeq \underline{\mathrm{End}}(L)$ déduit de b).
Enfin, si $J$ est comme dans c), le noyau est aussi égal à $J$. Et pour
$\mathcal{A}$ fixé, la corr.\ $L \mapsto u$ est une bijection entre
l'ensemble des $L$ satisfaisant a), et \uncertain{l'ensemble} des
\uncertain{homomorphismes} d'Algèbres $\mathcal{A} \to \underline{O}_S$.
\note{les deux dernières lignes de ce paragraphe sont serrées en interligne et se lisent mal ; « \ldots{} $\mathcal{A} \to \underline{O}_S$ » est suivi de « DES $L$ » et d'une fin illisible soulignée.}

\emph{Dém.} a), c) \uncertain{Ecrivons} $L^{\perp}$ au moyen de bases $x, x'$
\add{Si $u$ est un endom de $V$ i.e.\ section de $\mathcal{M}$}, on a
\[
  \mathrm{Tr}\, u \, x' \otimes x = \langle ux, x' \rangle
\]
\note{la ligne d'entrée « Écrivons \ldots{} bases $x, x'$ » est lue sous toutes réserves ; l'insertion « si $u$ est un endom de $V$ \ldots{} » est écrite au-dessus.}
qui est nulle ssi $ux$ est \uncertain{proportionnel} à $x$ i.e.\ $u(L) \subset L$,
d'où a) $\Leftrightarrow$ b). Évidemment si $u : \mathcal{A} \to \underline{O}_S \simeq \underline{\mathrm{End}}(L)$
est l'hom correspondant, il est clair que son noyau $J$ est un ss-fibré ann.\ par
$J \cdot \underline{L} = 0$, donc b) $\Rightarrow$ c).
Dans \uncertain{pour} un tel $J$, on \uncertain{a} \add{(\uncertain{nécessairement})}
$\underline{O}_S \xrightarrow{\sim} \mathcal{A}/J$ car $\mathcal{A}/J$ $\underline{O}_S$-Alg.\
libre de rg~1) donc $\mathcal{A} \simeq J + \underline{O}_S$ et $L$ est stable sous
$\mathcal{A}$, donc b) $\Leftrightarrow$ c)
\note{la fin de la page est très cursive ; seules les formules et les implications sont sûres.}

\page{52}
De plus, cet argument montre aussi que $J$ est uniquement déterminé comme
$\mathrm{Ker}\bigl(\mathcal{A} \to \underline{\mathrm{End}}(L)\bigr)$.
De plus, sous les conditions équivalentes a) b) c)
\struck{\ill{} c'est dire \ill{} $\det \mid \underline{J} = 0$ donc $J$ \ill{}}
\struck{\ill{}} sont loc.\ \uncertain{nulles} que \uncertain{base} de $\underline{J}$,
\struck{donc est sur} \struck{de la forme $\underline{L} \otimes$}
de sorte que $L = \mathrm{Ker}\, u$, \uncertain{soit} $M = \mathrm{Im}\, u$, c'est
un ss-fibré vectoriel dans $V$ tel que \struck{$L^{\perp} \otimes M$}
$L^{\perp} \otimes M = \underline{O}_S u = \underline{J} \subset \mathcal{A}$, d'où d).
Réciproquement, si $M$ est un ss-fibré inv.\ de $V$ tel que
$\check{L} \otimes M \subset \mathcal{A}$, évidemment $\check{L} \otimes M$ est un
ss-fibré inv.\ de $\mathcal{A}$ annulant $L$, donc c'est $J$, ce qui montre
\uncertain{à la fois} que $M$ est uniquement déterminé comme $\mathrm{Im}\, u$
($u$ base locale de $J$).

Remarque : prouvons d) $\Rightarrow$ b), \uncertain{car} \ill{} implique
$\mathcal{A} = \underline{O}_S + J$ (si on pose $J = L^{\perp} \otimes M$)
et $J$ étant nul sur $\underline{L}$ à fortiori l'invarie.

cqfd.

\emph{Corollaire.} Se donner une $\mathcal{A}$ (ss-fibré vectoriel de rg~2 de
$\mathcal{M}$ contenant \add{i.e.} $\underline{O}_S \cdot 1$ -- \uncertain{alors}
\uncertain{tel} $\mathcal{A}$ est \uncertain{une} ss-Alg.\ commutative !) ou se donner
$\mathcal{A}^{\perp} = \mathcal{H} \subset \widetilde{\mathfrak{g}}$, un sous-fibré
vectoriel de rang 2 \add{i.e.\ corang 1} de $\widetilde{\mathfrak{g}}$, revient au
même
\[
  (21) \qquad \mathrm{AlgSpec}(\mathcal{M}) \simeq \underbrace{\mathbb{P}(\widetilde{\mathfrak{g}})}_{P} \ \bigl(\simeq \check{\mathbb{P}}(\mathfrak{g})\bigr) \simeq \lbrace \ldots \rbrace
\]
\note{l'accolade finale de (21) contient : « solutions des (fibrés en droites) dans $P = \check{\mathbb{P}}(\widetilde{\mathfrak{g}}) = \mathbb{P}(\mathfrak{g})$ », lecture incertaine (peut-être « schéma des droites \ldots{} »). « Corollaire » est surchargé et souligné.}

\page{53}
(Il y a \struck{application} \add{un} morphisme)
\[
  D \longmapsto D \cap X = Y
\]
\[
  (22) \qquad \check{P} = \mathrm{Dr}(P) \xrightarrow{\ \sim\ } \mathrm{Div}_2^{+}(X)
\]
est un iso, qui associe : \struck{\ill{}} \add{à} $\mathcal{H} \subset \widetilde{\mathfrak{g}}$
le diviseur $Y \subset X$ des \struck{droites} \add{ss-fibrés inv.} de
$\widetilde{\mathfrak{g}}$ qui \emph{sont dans} $\mathcal{H}$, et satisfont
$q \mid \underline{L} = 0$.
\struck{C'est donc} Comme solution sur $S$, il s'exprime dans le
\uncertain{sous-schéma} de $\check{\mathbb{P}}(\mathcal{H}) = \check{\mathbb{P}}(\mathcal{H})$ défini par
l'annulation de $q$. On trouve donc, \struck{\ill{}} par \ill{} un iso.\ can.
\[
  (23) \qquad Y \simeq \operatorname{Spec} \mathcal{A}
\]
\note{la première égalité de la ligne « $\check{\mathbb{P}}(\mathcal{H}) = \check{\mathbb{P}}(\mathcal{H})$ » est ainsi sur la page ; l'un des deux $\mathbb{P}$ porte peut-être une autre marque.}

Dans l'iso précédent $X \simeq \mathbb{P}(V)$, \uncertain{via} $L \mapsto \underline{L}$,
on vient que $Y$ est le sous-schéma de $\mathbb{P}(V)$ formé des $L$ tels que
$\mathcal{A} L \subset L$ ou encore (si $\mathcal{A} = \underline{O}_S 1 + \underline{O}_S u$)
tels que $u L \subset L$, ou encore, si $u$ est inv.\ i.e.\ section de
$\check{\mathbb{V}}(\mathcal{M}^{*})$
\[
  (24) \qquad Y = X^{u}
\]
sous-schéma des \uncertain{invariants} sous $u$.
\note{l'accent et l'étoile de $\check{\mathbb{V}}(\mathcal{M}^{*})$ sont surchargés.}

En retournant réciproquement $\mathcal{A}$ pour un \uncertain{diviseur}
\uncertain{vérifiant} $Y$, on \ill{} que \struck{l'opération \ill{}}
\add{l'isomorphisme} (23) est fonctoriel, donc compatible avec
\uncertain{les} automorphismes de $u \in \Gamma(\mathcal{M}^{*})$ (i.e.\ ici
$u : V \to V$), on \uncertain{fait} opérer $u$ par \uncertain{automorphisme}
sur $\mathcal{M}^{*}$, et de $G$ de façon évidente sur $X$ \ldots

\page{54}
On trouve, en notant que $\mathcal{A}$ se déduit de l'action correspondante de
$G = \check{\mathbb{V}}(\mathcal{M}^{*})/\mathbb{G}_m$ :

\emph{Prop.} Pour que $g \in G(S)$ invarie $Y$, il faut et suffit qu'il invarie
$\mathcal{A}$ i.e.\ que \struck{\ill{}} 
\[
  \operatorname{int}(u)\, \mathcal{A} \ \bigl(\overset{\mathrm{def}}{=} u \mathcal{A} u^{-1}\bigr) = \mathcal{A}
\]
\note{la proposition est marquée d'un trait vertical dans la marge.}

Quand il en est ainsi, l'opération de $g$ sur $Y$ s'identifie via (23) à
l'op.\ de $g$ \add{(\ill{})} sur $\operatorname{Spec}(\mathcal{A})$ -- donc elle est
triviale ssi $g$ opère trivialement sur $\mathcal{A}$ i.e.\ ssi $u \in \Gamma(S, \mathcal{A})$.

Bien entendu, ceci caractérise $\mathcal{A}$ \uncertain{comme} engendré comme
sous-fibré vectoriel par les sections locales de $\mathcal{A}^{*}$ \ldots).

À \struck{\ill{}} $\mathcal{A}$, associons le ss-groupe
\[
  (25) \qquad \mathcal{T} = \check{\mathbb{V}}(\mathcal{A})^{*}/\mathbb{G}_m \subset G
\]
\note{la lettre de (25) est une majuscule ronde, lue $\mathcal{T}$ ; le quotient $\check{\mathbb{V}}(\mathcal{A})^{*}/\mathbb{G}_m$ est écrit sous elle, relié par un signe d'égalité vertical.}
c'est un sous-groupe \emph{commutatif} \emph{lisse} \emph{connexe} de dimension
\emph{relative}~1. Bien entendu, la correspondance $\mathcal{A} \mapsto \mathcal{T}$
est injective, $\mathcal{A}$ étant engendré par \struck{\ill{}} les sections de
$\check{\mathbb{V}}(\mathcal{A})^{*}$, qui est l'image inverse de $\mathcal{T}$ par
l'hom can \add{si 2 est inv.\ sur $S$} $\check{\mathbb{V}}(\mathcal{M})^{*} \to G$.
D'autre part, \struck{donc} \add{$\mathcal{T}$} s'y voit les propriétés
\uncertain{communes} (sauf la connexité \ldots)
\struck{-- \uncertain{Signalons} les ss-groupes \emph{jacobiens} de $G$
\uncertain{provenant} de tels $\mathcal{A}$} -- i.e.\ des ss-fibrés inv.\ \ldots
\note{la fin de la page, avec la boucle qui déplace « donc » et l'encadré « (sauf la connexité\ldots) », est surchargée ; la lecture de ces dernières lignes est incertaine.}

\page{55}
\add{\ldots} sont de la forme $\check{\mathbb{V}}(\mathcal{A}^{*})/\mathbb{G}_m$ -- i.e.\
sont des sous-groupes « \emph{jacobiens} ».
\note{cette ligne, qui achève la dernière phrase de la page 54, est écrite serrée en haut du feuillet.}
\struck{de $\mathfrak{g} = \mathrm{Lie}\, G$ + \ill{}}
\struck{Soit \ill{}}, considérons $\mathfrak{h} = \mathrm{Lie}\, \mathcal{T}$,
qui est un tel ss-fibré, ($\mathcal{T}$ lisse de dim.\ rel.~1) considérons le
\uncertain{sous}-groupe $\widetilde{\mathcal{T}}$ \struck{qu'il définit}
jacobien qu'il définit, qui est donc aussi le centralisateur de $\mathfrak{h}$ dans $G$.
Comme $\mathcal{T}$ est \emph{commutatif} il centralise $\mathfrak{h} = \mathrm{Lie}\, \mathcal{T}$,
donc $\mathcal{T} \subset \widetilde{\mathcal{T}}$, d'où égalité puisque
$\widetilde{\mathcal{T}}$ est connexe et $\mathcal{T}, \widetilde{\mathcal{T}}$ sont
lisses de dimension rel.~1 \ldots
\marginal{\ill{} \struck{il faut} cor.\ \uncertain{p.~52}}
\note{note marginale à gauche, écrite en oblique et soulignée de deux traits ; en face d'un trait vertical qui tient la phrase « donc aussi le centralisateur \ldots{} » ; « il faut » y est biffé, et le renvoi final (lu « p.~52 ») est incertain. Le reste de la page est blanc.}

\section*{Dictionnaire lié à $\mathrm{GP}(1) \simeq \mathrm{SO}(3)$}
\note{inscrit en haut à droite du feuillet de garde rose (page 56), qui ne porte rien d'autre ; ces mots sont de sa main et deviennent l'intitulé de la suite.}

\page{57}
$S$ schéma. On va définir des équivalences entre les \struck{\ill{}}
\add{\uncertain{neuf}} catégories suivantes \uncertain{attachées} à $S$ -- et pour $S$
variable, on aura des équivalences mutuelles de catégories fibrées sur
$\mathrm{Sch}_{/S}$ (en fait, des gerbes fppf sur $\mathrm{Sch}_{/S}$),
formant un « système transitif d'équivalences à isom.\ \uncertain{compris} » :

$C_{0_S} = \mathrm{Tors}\bigl(\mathrm{GP}(1)_S\bigr)$ catégorie des torseurs
\add{à droite} \struck{\ill{}} sous $\mathrm{GP}(1)_S$

$C_{1_S} = \mathrm{Fo}\bigl(\mathrm{GP}(1)_S\bigr)$ catégorie des « formes »
$G$ (\uncertain{fppf} \ldots{} \uncertain{f(ét)}) du \uncertain{GP(1)}, schéma en groupes $\mathrm{GP}(1)$

$C'_{1_S} = \mathrm{Fo}\bigl(\mathrm{SL}(2)_S\bigr)$ catégorie des formes $\widetilde{G}$ de $\mathrm{SL}(2)_S$

\struck{$C''_{1_S} = \mathrm{Fo}\bigl(\mathrm{GL}(2)_S\bigr)$ catégorie des formes \ill{}}

$C_{2_S} = \mathrm{Fo}\bigl(\mathfrak{gp}(1)_S\bigr)$ (*) catégorie des formes $\mathfrak{g}$ de l'Algèbre de Lie
$\mathfrak{gp}(1)_S$ de $\mathrm{GP}(1)_S$

$C'_{2_S} = \mathrm{Fo}\bigl(\mathfrak{sl}(2)_S\bigr)$ (*) \struck{catégorie des formes} \ldots{} $\widetilde{\mathfrak{g}}$ de
$\mathfrak{sl}(2)_S \simeq \mathrm{Lie}\, \mathrm{SL}(2)_S$

\struck{$C''_{2_S} = \mathrm{Fo}\bigl(\mathfrak{gl}(2)_S\bigr)$ (**) catégorie des formes \ldots{} $\mathfrak{gl}(2)_S$ \ldots}

$C_{3_S} = \mathrm{Quat}(\underline{O}_S)$ catégorie des Algèbres d'Azumaya \add{$\mathcal{M}$} de
rang~4 sur $S$ (ou Algèbres de \uncertain{Quaternions})

$C_{5_S} = \mathrm{Drpr}(S)$ catégorie des fibrés en droites projectives $X$ sur $S$

$C'_{5_S} = \mathrm{Quad}(S)$ catégorie des couples $(P, X)$ d'un fibré en plans projectifs
sur $S$ et d'une quadrique lisse $X \subset P$

$C_{4_S} = \mathrm{Modspquad}(S)$ catégorie des « Modules spéciaux quadratiques de rang~3 sur $S$ »
$(E, q, \omega)$, $E$ fibré vect.\ de rg~3, $\omega$ \ldots
\note{la liste est disposée en deux colonnes (le symbole, sa description) ; des accolades regroupent $C_1, C'_1, C''_1$, puis $C_2, C'_2, C''_2$, puis $C_5, C'_5$, et une flèche courbe relie $C_4$ à $C_3$. Les lignes $C''_1$ et $C''_2$ sont biffées ; les astérisques (*) et (**) sont dans la marge gauche. Le nom de $C_3$ surcharge un premier mot illisible ; celui de $C_4$ surcharge un premier nom (peut-être « $\mathrm{Spect}(S)^3$ »). Le dernier alinéa se poursuit à la page 58.}

\page{58}
Ici le signe (*), en marge indique que l'équivalence des 3 catégories en
question avec les \uncertain{6} autres n'est valable que si $S$ est premier à 2
-- en général on a seulement des foncteurs qui vont de chacune des
\uncertain{six} autres catégories vers chacune des trois marquées du signe (*),
mais ce ne sont pas des foncteurs pl.\ fidèles \ldots
\note{le nombre des « autres » catégories est mal formé ; il peut aussi se lire 8.}

Un dictionnaire complet devrait comprendre, \uncertain{avec} la description
explicite des $9 \times 8 = 72$ équivalences de catégories \uncertain{sus-}\ill{}
compter les \struck{\ill{}} $9 \cdot 8 \cdot 7 = 504$ isomorphismes
\uncertain{de} transitivité (qui \uncertain{seraient} être \uncertain{vérifiés}
\ldots) On fera un peu moins, mais \uncertain{on} essaiera de donner une image
claire de la situation.
\note{« $= 72$ » est inscrit au-dessus de $9 \times 8$ ; « $9 \cdot 8 \cdot 7 = 504$ » remplace un premier nombre biffé.}

\emph{Principe général.} On va définir des représentants typiques dans chaque
catégorie dans le cas où $S = S_0 = \operatorname{Spec} \mathbf{Z}$, et on
définira une opération de $G_0 = \mathrm{GP}(1)$ sur chacun.
On établira que ceci établit une isom.
\[
  \mathrm{GP}(1) \simeq \underline{\mathrm{Aut}}(\mathfrak{X}_0)
\]
de $\mathrm{GP}(1)$ avec le schéma des automorphismes de chacun des objets
$\mathfrak{X}$ \uncertain{envisagés} -- \ill{}
\note{la lettre gothique de $\underline{\mathrm{Aut}}(\mathfrak{X}_0)$ est lue $\mathfrak{X}$ ; l'indice de $G_0$ est surchargé d'une tache.}

\page{59}
bornant cependant à \struck{\ill{}} $S'_0 = S - \lbrace 2 \rbrace$ dans les trois
cas (*) où l'opération \add{sur $S_0$} est fidèle, mais l'hom.\ de groupes
\uncertain{associé} n'est pas épi.
Comme les 11 catégories \add{fibrées} envisagées sont des gerbes (fppf disons,
ou (ét)), il s'ensuit que chacune d'elles est can.\ équiv.\ à la cat.\ des
torseurs \add{à droite} sous le $\underline{\mathrm{Aut}}(\mathfrak{X}_0)$
correspondant, donc aussi à la cat.\ des torseurs à dr.\ sous $G_0$ (\ldots{} à
la \uncertain{précision} de \uncertain{sel} \ldots{} *). Ceci établira
l'\emph{existence} d'un multi-dictionnaire, mais on \uncertain{\ill{}} ne
\uncertain{dispensera} pas de le développer en détail.
\note{« 11 » est net sur la page, alors que la liste de la page 57, une fois biffées les lignes $C''_1$ et $C''_2$, en compte neuf ; la fin de la parenthèse « (\ldots{} *) » est illisible.}

\emph{I Représentants typiques des onze cat.\ fibrées.}

1) $C_{0_{S_0}}$ : le torseur trivial $P_0 = (G_{S_0})_d$ \uncertain{sous} $\mathrm{GP}(1)_{S_0} = G_{S_0}$

2) $C_{1_{S_0}}$ : la forme triviale $G_0 = G_{S_0} \overset{\mathrm{GP}(1)}{=} G_{S_0}$

\quad $C'_{1_{S_0}}$ \ldots{} $\widetilde{G}_0 = \mathrm{SL}(2)_0$ de $\mathrm{SL}(2)$

\struck{\quad $C''_{1_{S_0}}$ \ldots{} $= \mathrm{GL}(2)_0$ de $\mathrm{GL}(2)$}

3) $C_{2_{S_0}}$ \ldots{} $\mathfrak{g}_0 = \mathfrak{gp}(1)$ de $\mathfrak{gp}(1)$

\quad $C'_{2_{S_0}}$ \ldots{} $\widetilde{\mathfrak{g}}_0 = \mathfrak{sl}(2)$ de $\mathfrak{sl}(2)$

\struck{\quad $C''_{2_{S_0}}$ \ldots{} $\mathfrak{g}''_0 = \mathfrak{gl}(2)$ de $\mathfrak{gl}(2)$}

4) $C_{3_{S_0}}$ : l'Algèbre d'Azumaya triviale $\mathcal{M}_0 = \mathrm{M}_2(\underline{O}_{S_0})$

\struck{5) $C_{4_{S_0}}$ : $X_0 = \mathbb{P}^1_{S_0} = \mathbb{P}^1(\underline{O}^2_{S_0}) = \check{\mathbb{P}}(\underline{O}^2_{S_0})$}
\note{les lignes $C''_1$, $C''_2$ et le 5) sont biffés (le 5) de hachures obliques) ; sous le 5) biffé, un second indice « $C_{4_{S_0}}$ ». Les numéros 3) et 5) surchargent d'autres chiffres. Le torseur trivial porte la lettre $P_0$, lue sous un $\mathbb{R}$ ou $P$ surchargé ; l'indice « $d$ » est « à droite ».}

\page{60}
5) $C_{4_{S_0}}$ : $E_0 = \underline{O}^3_{S_0}$, $q_0(x, y, z) = xy + z^2$,
$\omega = e_1 \wedge e_2 \wedge e_3$ où $(e_1, e_2, e_3)$ est la base canonique
(\uncertain{\ill{}} $(x, y, z)$ pour $x e_1 + y e_2 + z e_3$)

6) $C_{5_{S_0}}$ : $X_0 = \mathbb{P}^1_{S_0} = \mathbb{P}(\underline{O}^2_{S_0}) = \check{\mathbb{P}}(\underline{O}^2_{S_0})$

\quad $C'_{5_{S_0}}$ : $P_0 = \mathbb{P}^2_{S_0} = \check{\mathbb{P}}(E_0)$, $C_0 = \mathbb{V}(q_0)$
(quadrique définie par la forme q.\ de 5°).
\note{$\mathbb{V}(q_0)$ : la lettre est la même que dans « $\mathrm{Var}(q)$ » de la page 46, notée ici par un $\mathbb{V}$ ajouré.}

\emph{II Opérations de $G_0$}

1) Translations : \emph{gauche} \ldots

\struck{2) $G_0$ opère sur $\mathfrak{g}_0 = \underline{\mathrm{Lie}}(G_0)$ par repr.\ adjointe
les op.\ de $\widetilde{G}_0$ et $\mathcal{L}_0$ sur leurs Algèbres de Lie respectives}
\note{ces deux lignes sont barrées et en partie encadrées ; le numéro, d'abord 3), est surchargé en 2).}

2) $G_0$ opère sur lui-même par automorphismes \uncertain{intérieurs} ;
les op.\ de $\widetilde{G}_0$ \struck{\ill{}} sur lui-même se factorisent par le quotient
$\widetilde{G}_0/\mathfrak{z}_0$ \struck{du \ill{}} \struck{en général} par le centre,
$\mathfrak{z}_0$ \ldots\ les \struck{\ill{}} iso can
\[
  (1) \qquad \widetilde{G}_0/\mathfrak{z}_0 = \mathrm{SL}(2)/\mu_2 \simeq \mathrm{GP}(1) = G_0
\]
\struck{$\mathrm{GL}(2)/\mathbb{G}_m \simeq \mathrm{GP}(1) = G_0$}
d'où opérations de $G_0$ sur $\widetilde{G}_0$ \struck{\ill{}}.
\note{le centre est noté par un $\mathfrak{z}$ en forme de 3 (cf.\ hand.md, § 4) ; l'ordre des mots de la ligne qui le précède est incertain.}

3) \struck{\ill{}} Les op.\ de $G_0$ sur $G_0$ et $\widetilde{G}_0$ définissent des op.\ sur
$\mathfrak{g}_0 = \underline{\mathrm{Lie}}(G_0)$ et $\widetilde{\mathfrak{g}}_0 = \mathrm{Lie}(\widetilde{G}_0)$.

4) L'opération \uncertain{régulière} de $\check{\mathbb{V}}(\mathrm{M}_0)^{*} \simeq \mathrm{GL}(2)$
\add{sur} $\mathrm{M}_0$ est \struck{\ill{}} \add{triviale} sur le centre $\simeq \mathbb{G}_m \subset \mathrm{GL}(2)$,
donc passe au quotient, en une opération \add{de} (2)
$\mathrm{GL}(2)/\mathbb{G}_m \overset{\mathrm{déf}}{=} \mathrm{GP}(1)$ sur $\mathrm{M}_0$.
\note{« sur » est inscrit sous le $\simeq$ de la ligne 4) ; « triviale » remplace un mot biffé. La page s'arrête ici.}

\end{document}
